\chapter{นวัตกรรมถ่านอัดแท่งเสริมประสิทธิภาพความร้อนด้วยคาร์บอนแบล็ค}
\label{chap:ch06}

\section*{แผนบริหารการสอนประจำบทที่ 6}
\addcontentsline{toc}{section}{แผนบริหารการสอนประจำบทที่ 6}

\noindent\textbf{หัวข้อเนื้อหาประจำบท}
\begin{enumerate}
    \item ความเป็นมาและแนวคิดการประยุกต์ใช้คาร์บอนแบล็คในพลังงานชุมชน
    \item ผลกระทบของสัดส่วนคาร์บอนแบล็คต่อค่าความร้อนและการเผาไหม้
    \item ความแข็งแรงเชิงกลและความหนาแน่นของถ่านอัดแท่ง
\end{enumerate}

\noindent\textbf{วัตถุประสงค์เชิงพฤติกรรม}
\begin{enumerate}
    \item อธิบายหลักการ ทฤษฎีฟิสิกส์ และแบบจำลองทางสิ่งแวดล้อมประจำบทที่ 6 ได้อย่างถูกต้อง
    \item วิเคราะห์ สังเคราะห์ และประยุกต์ใช้เครื่องมือตรวจวัดหรือกระบวนการแปรรูปพลังงานได้อย่างมีประสิทธิภาพ
    \item คำนวณ ประเมินผลทางเทอร์โมไดนามิกส์ ทัศนศาสตร์ หรือการลดก๊าซเรือนกระจกได้อย่างถูกต้องตามหลักวิชาการ
\end{enumerate}

\noindent\textbf{การวัดและประเมินผล}
\begin{table}[H]
\centering\small
\begin{tabularx}{\textwidth}{p{0.55\textwidth} p{0.25\textwidth} c}
\toprule
\textbf{เกณฑ์การประเมินผล} & \textbf{เครื่องมือวัดผล} & \textbf{สัดส่วนคะแนน} \\
\midrule
1. ทักษะการปฏิบัติการทดลองและการคำนวณ & แบบประเมินทักษะห้องแล็บ & 40\% \\
2. รายงานผลการทดลองและการวิเคราะห์ & การตรวจดาต้าชีทและรายงาน & 30\% \\
3. แบบฝึกหัดทบทวนและคำถามท้ายบท & แบบทดสอบท้ายบทเรียน & 20\% \\
4. จิตพิสัยและการมีส่วนร่วม & การเข้าเรียนและการตรงต่อเวลา & 10\% \\
\midrule
\textbf{รวมทั้งสิ้น} & & \textbf{100\%} \\
\bottomrule
\end{tabularx}
\end{table}
\newpage

% ==========================================================
% เนื้อหาบทเรียนประจำบทที่ 6
% ==========================================================

\section{ความเป็นมาและแนวคิดการประยุกต์ใช้คาร์บอนแบล็คในพลังงานชุมชน}
\index{ความเป็นมาและแนวคิดการประยุกต์ใช้คาร์บอนแบล็คในพลังงานชุมชน}

การศึกษาและทำความเข้าใจเกี่ยวกับความเป็นมาและแนวคิดการประยุกต์ใช้คาร์บอนแบล็คในพลังงานชุมชน (Origin and Concept of Carbon Black Briquette Research) มีความสำคัญอย่างยิ่งในงานฟิสิกส์สิ่งแวดล้อมและพลังงาน โดยมีรายละเอียดสาระสำคัญดังนี้

การผลิตถ่านอัดแท่งจากชีวมวลทางการเกษตร เช่น กะลามะพร้าว ซังข้าวโพด เหง้ามันสำปะหลัง และกิ่งไม้ตัดแต่ง มักประสบปัญหาข้อจำกัดด้านคุณภาพ ได้แก่ ค่าความร้อนไม่สม่ำเสมอ (เฉลี่ย 3,800-4,500 cal/g), อัตราการมอดไหม้เร็ว, เกิดควันและกลิ่นรบกวน, และความแข็งแรงเชิงกลต่ำแตกหักง่ายระหว่างการขนส่ง

คณะผู้วิจัยมหาวิทยาลัยราชภัฏรำไพพรรณี โดย ผศ.ดร.ชีวะ ทัศนา และอาจารย์นุชนาถ หนูเล็ก ได้ริเริ่มโครงการวิจัย "การเพิ่มคุณสมบัติทางความร้อนของถ่านอัดแท่งด้วยคาร์บอนแบล็ค" (Enhancement in Thermal Properties of Briquettes using Carbon Black)

คาร์บอนแบล็ค (Carbon Black) เป็นอนุภาคคาร์บอนบริสุทธิ์ขนาดนาโนเมตรที่ได้จากกระบวนการเผาไหม้หรือการสลายตัวด้วยความร้อนของสารไฮโดรคาร์บอน มีปริมาณคาร์บอนคงตัว (Fixed Carbon) สูงมากกว่า 95\% และแทบไม่มีความชื้นหรือสารระเหย การนำคาร์บอนแบล็คมาผสมกับผงถ่านชีวมวลในสัดส่วนที่เหมาะสมจึงเป็นนวัตกรรมที่ช่วยยกระดับคุณสมบัติทางพลังงานได้อย่างก้าวกระโดด

\begin{flamebox}[นวัตกรรมพลังงานความร้อน มรภ.รำไพพรรณี]
การเสริมประสิทธิภาพความร้อนของถ่านอัดแท่งด้วยคาร์บอนแบล็คช่วยยกระดับค่าความร้อนสูงเกิน 6,200 cal/g ลดควันดำ และเพิ่มความแข็งแรงเชิงกล ตอบโจทย์การใช้งานเชิงพาณิชย์และสิ่งแวดล้อม
\end{flamebox}

\section{ผลกระทบของสัดส่วนคาร์บอนแบล็คต่อค่าความร้อนและการเผาไหม้}
\index{ผลกระทบของสัดส่วนคาร์บอนแบล็คต่อค่าความร้อนและการเผาไหม้}

การศึกษาและทำความเข้าใจเกี่ยวกับผลกระทบของสัดส่วนคาร์บอนแบล็คต่อค่าความร้อนและการเผาไหม้ (Effects of Carbon Black Ratio on Thermal Performance) มีความสำคัญอย่างยิ่งในงานฟิสิกส์สิ่งแวดล้อมและพลังงาน โดยมีรายละเอียดสาระสำคัญดังนี้

จากการทดสอบในห้องปฏิบัติการฟิสิกส์พลังงาน มรภ.รำไพพรรณี โดยแปรสัดส่วนการผสมคาร์บอนแบล็คที่ 0\%, 5\%, 10\%, 15\%, และ 20\% โดยน้ำหนัก ร่วมกับแป้งมันสำปะหลังเป็นตัวประสาน (Binder 5\%):

ผลการทดลองเชิงประจักษ์:
1. การเพิ่มขึ้นของค่าความร้อน (Heating Value): ถ่านอัดแท่งชีวมวลมาตรฐาน (0\% Carbon Black) มีค่าความร้อน 4,450 cal/g เมื่อผสมคาร์บอนแบล็ค 10\% ค่าความร้อนเพิ่มขึ้นเป็น 5,380 cal/g และที่สัดส่วน 20\% ค่าความร้อนพุ่งสูงถึง 6,240 cal/g (เพิ่มขึ้นมากกว่า 40\%)
2. การลดลงของควันและกลิ่น (Smoke and Odor Reduction): เนื่องจากคาร์บอนแบล็คมีสารระเหยง่ายต่ำมาก การเผาไหม้จึงเป็นไปอย่างสมบูรณ์ ไร้ควันดำ และไม่ส่งกลิ่นฉุนรบกวน เหมาะสำหรับการประกอบอาหารและปิ้งย่างในอาคาร
3. อัตราการมอดไหม้ที่ยาวนาน (Prolonged Burn Time): โครงสร้างอสัณฐานที่แน่นหนาของคาร์บอนแบล็คช่วยชะลออัตราการสิ้นเปลืองเชื้อเพลิง ทำให้ถ่านอัดแท่งติดไฟทนนานกว่าถ่านชีวมวลทั่วไป 1.5 ถึง 2.0 เท่า

\begin{flamebox}[นวัตกรรมพลังงานความร้อน มรภ.รำไพพรรณี]
การเสริมประสิทธิภาพความร้อนของถ่านอัดแท่งด้วยคาร์บอนแบล็คช่วยยกระดับค่าความร้อนสูงเกิน 6,200 cal/g ลดควันดำ และเพิ่มความแข็งแรงเชิงกล ตอบโจทย์การใช้งานเชิงพาณิชย์และสิ่งแวดล้อม
\end{flamebox}

\section{ความแข็งแรงเชิงกลและความหนาแน่นของถ่านอัดแท่ง}
\index{ความแข็งแรงเชิงกลและความหนาแน่นของถ่านอัดแท่ง}

การศึกษาและทำความเข้าใจเกี่ยวกับความแข็งแรงเชิงกลและความหนาแน่นของถ่านอัดแท่ง (Mechanical Compressive Strength and Density) มีความสำคัญอย่างยิ่งในงานฟิสิกส์สิ่งแวดล้อมและพลังงาน โดยมีรายละเอียดสาระสำคัญดังนี้

ความแข็งแรงเชิงกลเป็นปัจจัยสำคัญในการป้องกันการแตกหักระหว่างการขนส่งและการเก็บกอง:
- ความหนาแน่นรวม (Bulk Density): เพิ่มขึ้นจาก 0.92 g/cm$^3$ เป็น 1.18 g/cm$^3$ ที่สัดส่วนคาร์บอนแบล็ค 15\%
- ความต้านทานแรงกดแตก (Compressive Strength): ทดสอบด้วยเครื่อง Universal Testing Machine (UTM) พบว่าสามารถรับแรงกดในแนวแกนได้มากกว่า 12.5 MPa ผ่านเกณฑ์มาตรฐานผลิตภัณฑ์อุตสาหกรรม (มอก. 2272-2549) อย่างดีเยี่ยม

\section{ขั้นตอนและวิธีปฏิบัติการทดลอง}
\index{ขั้นตอนการทดลองประจำบทที่ 6}

\subsection{สูตรและกระบวนการผลิตถ่านอัดแท่งผสมคาร์บอนแบล็คระดับห้องปฏิบัติการ}
\begin{conceptbox}[สูตรและกระบวนการผลิตถ่านอัดแท่งผสมคาร์บอนแบล็คระดับห้องปฏิบัติการ]
1. บดผงถ่านชีวมวล (กะลามะพร้าวหรือเศษไม้) และร่อนผ่านตะแกรงขนาด 60 เมช
2. ชั่งผงถ่านชีวมวล 800 กรัม ผสมกับผงคาร์บอนแบล็ค 150 กรัม (15\% w/w) ให้เป็นเนื้อเดียวกัน
3. เตรียมกาวแป้งมันสำปะหลัง: ละลายแป้งมัน 50 กรัม ในน้ำร้อน 300 มล. กวนจนเป็นเจลใสเหนียว
4. เทกาวแป้งลงในส่วนผสมคาร์บอน คลุกเคล้าด้วยเครื่องผสมเชิงกลเป็นเวลา 15 นาที
5. ป้อนเข้าเครื่องอัดแท่งแบบเกลียวอัด (Screw extruder) ที่อุณหภูมิกระบอกอัด 80-100$^{\circ}$C
6. นำแท่งถ่านที่ได้ไปตากแดด 2 วัน แล้วอบไล่ความชื้นในตู้อบที่ 80$^{\circ}$C เป็นเวลา 12 ชั่วโมง
\end{conceptbox}

\subsection{การทดสอบอัตราการสิ้นเปลืองเชื้อเพลิงและควันด้วยการต้มน้ำเดือด}
\begin{conceptbox}[การทดสอบอัตราการสิ้นเปลืองเชื้อเพลิงและควันด้วยการต้มน้ำเดือด]
1. วางเตาอั้งโล่มาตรฐาน ชั่งน้ำหนักถ่านอัดแท่งเริ่มต้น 500 กรัม บรรจุลงในเตา
2. จุดไฟด้วยหัวพ่นก๊าซเป็นเวลา 3 นาที บันทึกเวลาเริ่มติดไฟแดงสม่ำเสมอ
3. ตั้งหม้อต้มน้ำบรรจุน้ำกลั่น 3,000 มิลลิลิตร วัดอุณหภูมิน้ำเริ่มต้น
4. บันทึกเวลาที่น้ำเดือดถึง 100$^{\circ}$C และชั่งน้ำหนักน้ำที่ระเหยไป
5. ตรวจวัดความเข้มข้นของก๊าซ CO และควันด้วยเครื่องตรวจวัดไอเสีย Testo 350
6. คำนวณประสิทธิภาพความร้อนเชิงความร้อน (Thermal Efficiency, $\eta$)
\end{conceptbox}

\subsection{การทดสอบความต้านทานแรงกดแตก}
\begin{conceptbox}[การทดสอบความต้านทานแรงกดแตก]
1. ตัดแท่งถ่านให้มีขนาดความยาวเท่ากับเส้นผ่านศูนย์กลาง (อัตราส่วน 1:1) ขัดหน้าตัดให้เรียบขนาน
2. วางตัวอย่างถ่านบนแท่นกดของเครื่องทดสอบแรงกดสากล (Universal Testing Machine)
3. เพิ่มแรงกดในแนวดิ่งด้วยอัตราเร็วคงที่ 2 มม./นาที จนกระทั่งแท่งถ่านเกิดการแตกหัก
4. บันทึกค่าแรงกดสูงสุด ($F_{\max}$, N) และคำนวณค่าความเค้นกดแตก ($\sigma = F_{\max} / A$, MPa)
\end{conceptbox}

\section{ตารางบันทึกผลการทดลองและการอภิปรายผล}
\begin{table}[H]
\centering\small
\caption{ตารางบันทึกผลการทดลองและการตรวจวัดพารามิเตอร์ประจำบทที่ 6}
\label{tab:ch06_datasheet}
\begin{tabularx}{\textwidth}{p{0.3\textwidth} p{0.3\textwidth} X}
\toprule
\textbf{พารามิเตอร์ที่ตรวจวัด} & \textbf{ค่าที่บันทึกได้} & \textbf{การแปลผลและการวิเคราะห์ทางฟิสิกส์} \\
\midrule
การทดสอบชุดที่ 1 & ค่าอยู่ในเกณฑ์มาตรฐาน & สอดคล้องกับแบบจำลองทฤษฎี \\
การทดสอบชุดที่ 2 & ค่ามีแนวโน้มเพิ่มขึ้นตามสัดส่วน & ยืนยันสมมติฐานการวิจัยเชิงประจักษ์ \\
ประสิทธิภาพรวม & ร้อยละ 88.5 & แสดงผลสัมฤทธิ์ในระดับยอมรับได้ \\
\bottomrule
\end{tabularx}
\end{table}

\section{คำถามทบทวนและแบบฝึกหัดท้ายบท}
\begin{enumerate}
    \item จงอธิบายกลไกทางฟิสิกส์และเคมีที่ทำให้การเติมคาร์บอนแบล็คช่วยเพิ่มค่าความร้อนของถ่านอัดแท่งได้อย่างมีนัยสำคัญ
    \item หากผสมคาร์บอนแบล็คในสัดส่วนสูงเกิน 25\% โดยน้ำหนัก จะส่งผลเสียต่อกระบวนการขึ้นรูปและคุณสมบัติของแท่งถ่านอย่างไร
    \item จงเปรียบเทียบผลการทดสอบ Water Boiling Test ระหว่างถ่านไม้ธรรมดากับถ่านอัดแท่งเสริมคาร์บอนแบล็ค 15\%
\end{enumerate}
